
Dette prosjektet ble gjennomført i samarbeid med Kystverket. Målet med prosjektet var å utforske og evaluere alternative visualiseringsteknikker for store datasett, samt å teste om slike visualiseringsmetoder kunne gi verdi for kunden.

Systemet ble utviklet som en applikasjon med en hjemmeside. React og TypeScript ble brukt for å lage applikasjonen, med MapLibre GL JS som basiskart og Deck.gl som datavisualiseringslag. Backend ble levert av kunden, og REST API-endepunkter ble brukt for å hente data som frontend visualiserer.

Flere visualiseringsmetoder ble implementert, inkludert klynging, varmekart, signalstyrkevisualisering og visualisering av anomaligrupper. Disse ble brukt til å evaluere hvor effektivt slike metoder kan forenkle store datasett til visuelt forståelige representasjoner.

Resultatene viser at kombinasjonen av flere ulike visualiseringsmetoder kan gi bedre innsikt for maritime operatører. Gruppen fikk positive tilbakemeldinger på den implementerte løsningen. Dette viser at prototyping med open source-verktøy og testdata kan gi verdi før full produksjon av slik funksjonalitet blir laget for interne verktøy, slik at man kan evaluere om funksjonaliteten gir verdi før produksjon. 


